% Laboratorio 3 - informe técnico.
% Compilar desde esta carpeta: pdflatex informe.tex (dos veces para el índice).
% Resultados de 51 ejecuciones reales; gráficas reproducibles desde los CSV.
\documentclass[11pt,letterpaper]{article}
\usepackage[utf8]{inputenc}
\usepackage[T1]{fontenc}
\usepackage{lmodern}
\usepackage{geometry}
\usepackage{xcolor}
\usepackage{graphicx}
\usepackage{booktabs}
\usepackage{array}
\newcolumntype{L}[1]{>{\raggedright\arraybackslash}p{#1}}
\usepackage{amsmath}
\usepackage{hyperref}

\geometry{left=2.7cm,right=2.7cm,top=2.6cm,bottom=2.7cm}
\definecolor{azul}{HTML}{10385F}
\definecolor{azulclaro}{HTML}{EAF1F6}
\definecolor{gris}{HTML}{526171}
\hypersetup{colorlinks=true,linkcolor=azul,urlcolor=azul,pdftitle={Laboratorio 3 - MPI},pdfauthor={Daniel Barillas}}
\setlength{\parindent}{0pt}
\setlength{\parskip}{0.55em}
\setlength{\emergencystretch}{2em}
\raggedbottom
\clubpenalty=10000
\widowpenalty=10000
\renewcommand{\figurename}{Figura}
\renewcommand{\contentsname}{Contenido}
\renewcommand{\abstractname}{Resumen}
\renewcommand{\refname}{Referencias}
\pagestyle{myheadings}
\markright{\small Laboratorio 3 \quad | \quad Computación Paralela y Distribuida}

\newcommand{\heading}[1]{\section{\textcolor{azul}{#1}}}
\newcommand{\subheading}[1]{\subsection{\textcolor{azul}{#1}}}
\newcommand{\evidence}[2]{%
  \begin{center}
  \includegraphics[width=\linewidth,height=0.38\textheight,keepaspectratio]{../evidencias/#1.png}\\[0.3em]
  \refstepcounter{figure}\label{fig:#1}
  {\small\textcolor{gris}{\textbf{Figura \thefigure.} #2}}
  \end{center}}
\newcommand{\chart}[2]{%
  \begin{center}
  \includegraphics[width=0.96\linewidth]{figuras/#1.pdf}\\[0.3em]
  \refstepcounter{figure}\label{fig:#1}
  {\small\textcolor{gris}{\textbf{Figura \thefigure.} #2}}
  \end{center}}

\begin{document}
\thispagestyle{empty}
\vspace*{2.4cm}
{\color{azul}\rule{\linewidth}{1.8pt}}\\[1.1cm]
{\Large Universidad del Valle de Guatemala}\\[0.4cm]
{\large Computación Paralela y Distribuida}\\[2.2cm]
{\Huge\bfseries\color{azul} Laboratorio 3}\\[0.35cm]
{\Large Comunicación bloqueante y no bloqueante con MPI}\\[0.8cm]
{\color{azul}\rule{0.62\linewidth}{0.8pt}}\\[2.0cm]
\begin{tabular}{@{}ll@{}}
\textbf{Estudiante:} & Daniel Barillas \\
\textbf{Carné:} & 22193 \\
\textbf{Sección:} & 40 \\
\textbf{Docente (enunciado):} & Marlon Fuentes \\
\textbf{Entrega vigente:} & Sábado 3 de octubre de 2026, 23:59 \\
\end{tabular}
\vfill
{\small\color{gris}Informe experimental con ocho capturas originales,
cuatro gráficas y 51 mediciones verificadas.\\
Preparado el 2 de octubre de 2026.}
\newpage

\begin{abstract}
Este informe documenta el diseño de cuatro programas de comunicación
punto a punto con MPI. Los dos primeros emplean llamadas bloqueantes para
resolver un intercambio Ping-Pong y un anillo con un único token. Los dos
restantes estudian comunicación no bloqueante mediante una recepción
anticipada con trabajo aritmético y un productor-consumidor que transfiere
un arreglo por chunks. Se define un protocolo reproducible de medición con
\texttt{MPI\_Wtime}, validación de datos y repeticiones por configuración.
Se analizan 51 ejecuciones realizadas por el estudiante en WSL: tres
repeticiones de cada una de 17 configuraciones. El RTT mediano de Ping-Pong
aumenta de 0.495 a 50.171 microsegundos entre los mensajes mínimo y máximo.
El anillo con llamadas separadas registra 2.242 ms, frente a 3.290 ms con
Sendrecv; esta diferencia se interpreta considerando el tráfico adicional
de centinelas. La recepción anticipada realiza cálculo antes de confirmar
la llegada del mensaje, sin medir una aceleración frente a una referencia
bloqueante. El pipeline obtiene su menor tiempo mediano, 2.279 ms, con
chunks de 65\,536 enteros. Se conservan todas las observaciones y se limita
la interpretación al entorno y a los tamaños ensayados.
\end{abstract}

{\small\setlength{\parskip}{0pt}\tableofcontents}
\newpage

\heading{Introducción y alcance}

MPI (Message Passing Interface) permite que procesos independientes
intercambien datos mediante mensajes. Cada proceso conoce su identificador
mediante \texttt{MPI\_Comm\_rank} y el tamaño de su comunicador mediante
\texttt{MPI\_Comm\_size}. En los cuatro programas se utiliza
\texttt{MPI\_COMM\_WORLD}, que agrupa a los procesos iniciados por
\texttt{mpirun}. La coordinación correcta depende de emparejar origen,
destino, etiqueta y comunicador en cada transferencia.

El enunciado divide el trabajo en dos partes: comunicación con bloqueo y
comunicación sin bloqueo. Cada parte contiene dos ejercicios. La tabla de
evaluación del PDF muestra cuatro partes de 25\,\%, aunque solo las primeras
dos están nombradas en el texto; por ello este informe presenta cuatro
secciones de ejercicios sin asignarles oficialmente un porcentaje.

\begin{center}
\begin{tabular}{@{}L{0.20\linewidth}L{0.33\linewidth}L{0.38\linewidth}@{}}
\toprule
\textbf{Ejercicio} & \textbf{Comunicación} & \textbf{Evidencia requerida} \\
\midrule
Ping-Pong & \texttt{MPI\_Send} y \texttt{MPI\_Recv} & Tiempo frente a tamaño de mensaje \\
Token Ring & \texttt{MPI\_Sendrecv} y variante separada & Recorrido correcto y explicación de diferencias \\
Recepción anticipada & \texttt{MPI\_Irecv}, \texttt{MPI\_Isend}, \texttt{MPI\_Test} & Trabajo efectuado y tiempo frente a tamaño \\
Pipeline por chunks & \texttt{MPI\_Isend}, recepción por chunk y STOP & Tiempo total frente a tamaño de chunk \\
\bottomrule
\end{tabular}
\end{center}

\heading{Fundamentos de comunicación MPI}

\subheading{Llamadas bloqueantes}

En una recepción bloqueante, la llamada retorna cuando el buffer contiene
el mensaje recibido. En un envío bloqueante, el retorno indica que el
buffer del emisor puede reutilizarse; no implica por sí solo que el receptor
haya terminado de procesar los datos. La pareja de llamadas exige un orden
compatible entre procesos. Si todos los participantes esperan una
operación que nadie puede completar, se produce un bloqueo circular.

\texttt{MPI\_Sendrecv} combina en una sola llamada un envío y una recepción.
La operación es bloqueante, pero el subsistema MPI coordina ambos lados y
facilita desplazamientos en topologías como un anillo. Los buffers de envío
y recepción deben ser distintos. La alternativa con \texttt{MPI\_Send} y
\texttt{MPI\_Recv} separados requiere planear explícitamente el orden.

\subheading{Llamadas no bloqueantes}

\texttt{MPI\_Irecv} y \texttt{MPI\_Isend} crean una petición y retornan
antes de que necesariamente termine la transferencia. \texttt{MPI\_Test}
consulta una petición sin esperar hasta su finalización. Esto permite
alternar comunicación y cálculo dentro del mismo proceso. Una petición de
envío debe completarse antes de modificar o liberar su buffer.

\subheading{Medición temporal}

\texttt{MPI\_Wtime} entrega segundos transcurridos según el reloj local
del proceso. Cada duración se calcula restando dos lecturas tomadas en el
\emph{mismo} rank; no se restan marcas temporales de ranks distintos. Se
realiza una ronda de calentamiento fuera del cronómetro en el barrido de
Ping-Pong mediante la opción explícita \texttt{-{}-warmup}; sin esa opción,
el caso básico ejecuta exactamente N rondas, sin intercambios adicionales.
Se realizan tres repeticiones por configuración. La mediana reduce la
influencia de una corrida aislada sobre el resumen, pero no sustituye
la presentación de la dispersión.

\heading{Entorno y método experimental}

Las pruebas se realizaron localmente en Ubuntu bajo WSL con Open MPI. El
código utiliza C11 y se compila mediante \texttt{mpicc}, con optimización
\texttt{-O2} y advertencias \texttt{-Wall -Wextra -Wpedantic}. El entorno
disponible se consultó durante la preparación del informe, sin volver a
ejecutar los experimentos:

\begin{center}
\begin{tabular}{@{}L{0.28\linewidth}L{0.63\linewidth}@{}}
\toprule
\textbf{Componente} & \textbf{Entorno consultado} \\
\midrule
Sistema & Ubuntu 22.04.3 LTS sobre WSL2 \\
Kernel & \texttt{6.18.33.2-microsoft-standard-WSL2} \\
MPI & Open MPI 4.1.2, tanto \texttt{mpicc} como \texttt{mpirun} \\
Procesador & Intel Core i5-10210U, frecuencia nominal 1.60 GHz \\
Topología visible & 4 núcleos, 2 hilos por núcleo; 8 CPU lógicas \\
Generación de gráficas & Python 3.13.1 y Matplotlib 3.10.1 en Windows \\
\bottomrule
\end{tabular}
\end{center}

Esta consulta describe el entorno instalado, no una traza histórica de
carga, frecuencia o afinidad de cada corrida. El comando del anillo incluye
\texttt{-{}-oversubscribe}; esa opción permite superar los slots que informe
MPI, pero no demuestra que efectivamente se hayan excedido ni garantiza
un núcleo dedicado por rank.

Para cada experimento se mantienen constantes los parámetros de control
definidos en la tabla metodológica. En recepción anticipada se fija el
lote, no la cantidad total de cálculo. Los programas imprimen una línea CSV
al final, además de un resumen legible. Las capturas documentan la
ejecución y los CSV conservan los números exactos para las gráficas.
Las 51 filas de los CSV coinciden con la última línea de sus registros
originales locales. Hay 17 configuraciones completas, con repeticiones
1, 2 y 3, sin filas faltantes ni duplicadas. Las capturas corresponden a
demostraciones adicionales: sus tiempos no se mezclan con el barrido.
Las mediciones de Ping-Pong proceden de un protocolo con una ronda de
calentamiento fuera del cronómetro. En la versión final esa preparación
requiere \texttt{-{}-warmup}, preservando el protocolo del barrido sin
alterar los CSV históricos. Las dos capturas renovadas sí muestran la
versión final: el caso básico ejecuta N rondas exactas y el caso de mayor
tamaño activa explícitamente el calentamiento. Sus tiempos son demostrativos
y no se incorporan a las tablas del barrido.

Las tablas presentan mediana, media, desviación estándar muestral (DE)
y mínimo--máximo. La DE se calcula con divisor $n-1$, con $n=3$. No se
eliminó ninguna observación. Las figuras muestran las tres corridas, la
mediana y el rango observado; las barras \emph{no} son intervalos de
confianza. Las líneas entre medianas son guías de lectura, no modelos
ajustados. Los valores se redondean a tres decimales para presentación.

Los puntos de las tres corridas tienen un pequeño desplazamiento
horizontal para distinguirlos; no representan tamaños de mensaje distintos.

El análisis es descriptivo: con tres repeticiones por configuración no
se afirma significancia estadística ni una mejora generalizable. El
script \texttt{analizar\_resultados.py} valida los datos y reproduce las
tablas; \texttt{graficas\_resultados.py} produce las cuatro figuras sin
lanzar procesos MPI ni modificar los CSV.

\subheading{Presentación de consola y trazabilidad}

Los cuatro programas comparten una presentación textual en terminal:
encabezados azules, rutas de mensajes cian, comprobaciones verdes y notas
amarillas. Los diagramas explican el recorrido de los datos, pero no
representan una escala física de tiempo. Las barras de veinte celdas
indican finalización o visitas verificadas; la medición numérica siempre
aparece por separado. La variable \texttt{NO\_COLOR} elimina los códigos
ANSI cuando se preparan registros de texto. La última línea se conserva
como CSV puro para evitar transcribir tiempos desde las capturas.

La implementación separa la lógica MPI de la presentación mediante
\texttt{common.h} y \texttt{visual.h}. Cada programa documenta sus
argumentos, roles de los ranks, etiquetas, buffers, orden de llamadas,
verificación y límites de la región cronometrada. Los dibujos se imprimen
\emph{después} de calcular el tiempo con \texttt{MPI\_Wtime}; así no
alteran las comparaciones solicitadas.

\begin{center}
\begin{tabular}{@{}L{0.27\linewidth}L{0.30\linewidth}L{0.34\linewidth}@{}}
\toprule
\textbf{Experimento} & \textbf{Variable independiente} & \textbf{Control} \\
\midrule
Ping-Pong & 1, 64, 1024, 16384, 65536 enteros & 1000 rondas, 2 ranks \\
Token Ring & Dos variantes de comunicación & 5 ranks, 1000 vueltas \\
Recepción anticipada & Los mismos tamaños de mensaje & 1000 operaciones por consulta, 2 ranks \\
Pipeline & Chunks de 256 a 65536 enteros & 1\,048\,576 enteros totales, 2 ranks \\
\bottomrule
\end{tabular}
\end{center}

\clearpage
\heading{Parte 1: comunicación con bloqueo}

\subheading{Ping-Pong}

\subsubsection*{Inciso 1: entero y N intercambios}
\evidence{01_ping_pong_entero}{Versión final: 10 rondas exactas con un entero y calentamiento desactivado. El mensaje regresa sin cambios.}

El rank 0 genera un buffer de enteros y lo envía a rank 1. Este lo recibe y
lo devuelve sin modificarlo. Una ronda consta de un envío y una devolución;
el usuario fija N mediante el primer argumento. Para cumplir simultáneamente
el caso básico y el estudio de tamaños, el segundo argumento controla el
número de enteros: con valor 1 se envía un entero, y con valores mayores se
envía un arreglo de enteros. El último mensaje devuelto se compara con el
patrón original antes de informar que la prueba fue correcta.
Sin \texttt{-{}-warmup}, el programa final hace exactamente N intercambios;
con la opción añade una única ronda preparatoria, excluida de la medición.

La implementación separa las dos perspectivas en
\texttt{round\_as\_origin} y \texttt{round\_as\_echo}. La salida dibuja
PING y PONG con etiquetas diferentes, muestra N rondas completadas y
presenta el tiempo solo después de verificar el mensaje. Esto permite
relacionar la captura con las llamadas MPI del código.

\clearpage
\subsubsection*{Inciso 2: MPI\_Wtime, tamaños y tendencia}
\evidence{02_ping_pong_tamanos}{Versión final con \texttt{-{}-warmup}: una ronda preparatoria no medida y 1000 rondas medidas con 1024 enteros por mensaje.}

La duración total medida en rank 0 se divide entre N para obtener la
latencia media de ida y vuelta:
\[
  T_{\mathrm{RTT}}=\frac{t_{\mathrm{fin}}-t_{\mathrm{inicio}}}{N},
  \qquad
  T_{\mathrm{ida,estimada}}=\frac{T_{\mathrm{RTT}}}{2}.
\]
La segunda expresión es una estimación bajo simetría aproximada, no una
medición independiente de una dirección. La gráfica utiliza bytes
por mensaje en el eje horizontal y $T_{\mathrm{RTT}}$ en el vertical.

\subsubsection*{Resultados del barrido}
Las tres repeticiones de cada tamaño usan 1000 rondas y dos ranks.
Las 1000 rondas aportan un promedio por corrida; no representan 1000
muestras estadísticas independientes.

\begin{center}\small
\begin{tabular}{@{}rrrrr@{}}
\toprule
\textbf{Enteros} & \textbf{Bytes} & \textbf{Mediana ($\mu$s)} & \textbf{Media $\pm$ DE ($\mu$s)} & \textbf{Mín.--máx. ($\mu$s)}\\
\midrule
1 & 4 & 0.495 & $0.545\pm0.118$ & 0.461--0.680\\
64 & 256 & 0.578 & $0.579\pm0.013$ & 0.566--0.592\\
1024 & 4096 & 3.625 & $3.824\pm0.489$ & 3.467--4.381\\
16384 & 65536 & 13.846 & $13.717\pm0.316$ & 13.357--13.947\\
65536 & 262144 & 50.171 & $49.659\pm2.756$ & 46.683--52.124\\
\bottomrule
\end{tabular}\end{center}

\clearpage
\chart{01_ping_pong}{RTT frente a tamaño del mensaje. Ambos ejes son logarítmicos; se muestran las tres corridas y su rango observado.}

\subsubsection*{Interpretación}
El RTT mediano aumenta en los cinco tamaños ensayados. Entre 4 y 256 bytes
pasa de 0.495 a 0.578 $\mu$s; a 262\,144 bytes alcanza 50.171 $\mu$s.
La razón entre extremos es 101.36, aunque el tamaño del mensaje crece
65\,536 veces. Esto es compatible con un costo fijo por intercambio y
otro dependiente del tamaño, pero no se estimó cada componente por
separado ni se identificaron umbrales internos de protocolos MPI.

El mensaje de un entero muestra una media superior a su mediana debido
a la corrida de 0.680 $\mu$s. Se mantiene ese valor en la tabla y en la
figura. En el tamaño máximo, el rango 46.683--52.124 $\mu$s muestra que
el costo depende también de la variación entre ejecuciones. Los valores
no son mediciones de una red entre computadoras distintas.

Las capturas iniciales acreditan el caso básico y el mensaje mayor;
el análisis cuantitativo utiliza \texttt{ping\_pong.csv}. La validación
comprueba todo el último eco contra $\mathrm{valor}[i]=i\bmod1009$.

\clearpage
\subheading{Token Ring}

\subsubsection*{Inciso 1: anillo con más de cuatro procesos}
\evidence{03_token_ring_send_recv}{Anillo con envíos y recepciones separados.}

Se usan al menos cinco ranks. Para un rank $i$ de un comunicador de tamaño
$P$, el vecino anterior y el siguiente son
\[
  \mathrm{src}(i)=(i-1+P)\bmod P,\qquad
  \mathrm{dst}(i)=(i+1)\bmod P.
\]
Rank 0 comienza con un único token válido. Cada recepción incrementa el
valor, lo cual permite verificar que el token visitó todos los procesos.
Tras $V$ vueltas, el token devuelto a rank 0 debe valer $P\,V$, y cada rank
debe registrar exactamente $V$ visitas.

En la primera variante, rank 0 envía para iniciar cada vuelta y después
recibe el token de su vecino anterior. Los otros ranks reciben primero y
luego envían. Este orden evita que todos queden esperando dentro de un
envío bloqueante.

\clearpage
\subsubsection*{Incisos 2 y 3: Sendrecv y diferencia con Send/Recv}
\evidence{04_token_ring_sendrecv}{Anillo con MPI\_Sendrecv: evidencia para el ajuste del programa y su comparación con llamadas separadas.}

En la segunda variante, todos los ranks ejecutan
\texttt{MPI\_Sendrecv} en cada paso: el rank que no posee el token envía
un centinela $-1$, mientras el poseedor envía el valor válido. Una vuelta
consta de $P$ pasos; sigue existiendo un único token válido.

La diferencia central es que \texttt{MPI\_Sendrecv} combina las dos
operaciones bloqueantes en una llamada y coordina su progreso. Con
\texttt{MPI\_Send} y \texttt{MPI\_Recv} separados, el programador debe
establecer el orden seguro de cada rank. El centinela hace explícito el
estado vacío del anillo y permite que todos llamen a \texttt{MPI\_Sendrecv}
en cada paso, aunque solo un mensaje contenga el token verdadero.

El resumen visual dibuja la ruta $0\rightarrow1\rightarrow\cdots
\rightarrow P-1\rightarrow0$ y una barra de visitas por rank. Las
barras confirman el contador esperado, mientras el valor numérico final
detecta un salto omitido o duplicado.

\subsubsection*{Resultados y carga de mensajes}
El barrido utiliza cinco procesos y 1000 vueltas. Las seis filas registran
token final 5000; el programa comprueba además 1000 visitas por rank.
El tiempo local de rank 0 incluye la barrera final y excluye el
\texttt{MPI\_Gather} posterior que reúne las visitas.

\begin{center}\small
\begin{tabular}{@{}lrrr@{}}
\toprule
\textbf{Modo} & \textbf{Mediana (ms)} & \textbf{Media $\pm$ DE (ms)} & \textbf{Mín.--máx. (ms)}\\
\midrule
\texttt{send\_recv} & 2.242 & $2.245\pm0.122$ & 2.125--2.369\\
\texttt{sendrecv} & 3.290 & $3.314\pm0.211$ & 3.116--3.535\\
\bottomrule
\end{tabular}\end{center}

\clearpage
\chart{02_token_ring}{Comparación descriptiva del anillo. Las implementaciones conservan un token válido, pero generan diferente cantidad de mensajes.}

\subsubsection*{Diferencia funcional y diferencia temporal}
La mediana de Sendrecv es 1.467 veces la de las llamadas separadas,
equivalente a 46.71\,\% más tiempo. No es una comparación aislada de
primitivas equivalentes. Con llamadas separadas se transfieren cinco
mensajes por vuelta, es decir, 5000 mensajes de token. En Sendrecv hay
cinco pasos por vuelta y cinco envíos por paso: 25\,000 mensajes,
incluidos los centinelas. Se excluyen de estas cuentas las barreras
y la recopilación de visitas.

La variante Sendrecv simplifica la coordinación del desplazamiento,
pero su implementación efectúa más trabajo de comunicación. El resultado
no demuestra que la primitiva Sendrecv sea intrínsecamente más lenta ni
que pierda su utilidad para evitar dependencias cíclicas. Las dos
variantes cumplen la misma propiedad lógica del recorrido, no la misma
carga de mensajes. Las capturas de una vuelta muestran una visita por
rank y token final 5 en ambos modos.

\clearpage
\heading{Parte 2: comunicación sin bloqueo}

\subheading{Recepción anticipada}

\subsubsection*{Incisos 1 y 2: recepción, aritmética y MPI\_Test}
\evidence{05_recepcion_entero}{Recepción básica de un entero y trabajo aritmético.}

Rank 1 inicia \texttt{MPI\_Irecv} antes de hacer aritmética. A continuación
ejecuta un lote fijo de operaciones y llama a \texttt{MPI\_Test}; repite
esta secuencia hasta que la petición se completa. Rank 0 prepara el buffer,
inicia \texttt{MPI\_Isend} y espera antes de liberarlo. El caso de un
elemento cumple el intercambio de un entero y los arreglos permiten
estudiar mensajes mayores sin cambiar el protocolo.

El receptor registra el primer entero, la suma validada, el número total
de operaciones aritméticas y el tiempo transcurrido desde antes de publicar
la recepción hasta que \texttt{MPI\_Test} informa su finalización. También
se registra por separado el tiempo local de envío del rank 0. El tiempo
del receptor incluye cómputo y sondeo; por ello no debe interpretarse como
latencia pura de red. La cantidad de trabajo realizado antes de recibir
puede cambiar entre corridas según la planificación de procesos.

La línea de tiempo de la terminal distingue el envío de rank 0 de la
secuencia \texttt{Irecv} $\rightarrow$ cálculo $\rightarrow$
\texttt{MPI\_Test} de rank 1. \texttt{do\_arithmetic} mantiene visible
el estado del cálculo y \texttt{verify\_and\_sum} recorre todos los
enteros después de completar la petición.

\texttt{MPI\_Test} comprueba la finalización de la \emph{petición MPI},
no la operación aritmética. Esta distinción concreta el sondeo solicitado:
el cálculo se ejecuta sincrónicamente por lotes mientras la recepción sigue
pendiente, y su estado y contador se imprimen al confirmar el mensaje.

\clearpage
\subsubsection*{Inciso 3: medición según tamaño del mensaje}
\evidence{06_recepcion_tamanos}{Recepción anticipada con un mensaje mayor: datos, trabajo efectuado y tiempos locales.}

\subsubsection*{Resultados temporales}
En todas las corridas se utilizan lotes de 1000 iteraciones aritméticas.
El contador de operaciones corresponde a iteraciones del bucle, no a
FLOP ni a instrucciones de procesador medidas.

\begin{center}\small
\begin{tabular}{@{}rrrrr@{}}
\toprule
\textbf{Bytes} & \textbf{Mediana ($\mu$s)} & \textbf{Media $\pm$ DE ($\mu$s)} & \textbf{Mín.--máx. ($\mu$s)} & \textbf{Envío ($\mu$s)}\\
\midrule
4 & 6.701 & $6.935\pm0.680$ & 6.402--7.701 & 2.801\\
256 & 7.001 & $7.601\pm1.683$ & 6.301--9.502 & 3.601\\
4096 & 19.804 & $21.171\pm2.543$ & 19.604--24.105 & 18.703\\
65536 & 48.710 & $49.442\pm2.289$ & 47.609--52.008 & 47.409\\
262144 & 121.223 & $124.056\pm19.213$ & 106.417--144.528 & 118.823\\
\bottomrule
\end{tabular}\end{center}
{\small La última columna es la mediana del tiempo local de envío de rank 0;
las demás columnas describen recepción, cálculo y sondeos de rank 1.}

\clearpage
\chart{03_recepcion_anticipada}{Tiempo del receptor frente a tamaño de mensaje. El eje horizontal es logarítmico y el vertical es lineal.}

\subsubsection*{Trabajo realizado y datos recibidos}
\begin{center}
\begin{tabular}{@{}rrr@{}}
\toprule
\textbf{Enteros} & \textbf{Iteraciones mín.--máx.} & \textbf{Suma verificada}\\
\midrule
1 & 1000--2000 & 0\\
64 & 1000--1000 & 2016\\
1024 & 2000--3000 & 508641\\
16384 & 2000--3000 & 8165256\\
65536 & 1000--5000 & 33006624\\
\bottomrule
\end{tabular}\end{center}

La mediana crece de 6.701 a 121.223 $\mu$s, un factor de 18.09. El
trabajo total no es constante ni crece monótonamente con el mensaje:
a 65\,536 enteros se realizan entre 1000 y 5000 iteraciones. Esa variación
depende de cuándo el sondeo confirma la petición y no mide por sí sola
el rendimiento del transporte. El bucle ejecuta al menos un lote antes
de la primera consulta, incluso si el mensaje pudiera completarse antes.

Los tiempos de envío y recepción tienen intervalos locales distintos;
restarlos no aísla el costo del cálculo. La evidencia demuestra trabajo
entre la publicación y la confirmación de la recepción. Sin una variante
bloqueante equivalente y con el mismo trabajo total, no se puede
cuantificar una aceleración por solapamiento. Las 15 sumas coinciden con
el patrón $i\bmod1009$, y el estado aritmético visible impide que el
resultado del cálculo sea descartado por el compilador.

\clearpage
\subheading{Productor-consumidor: pipeline por chunks}

\subsubsection*{Incisos 1 y 2: chunks iguales y STOP con Isend}
\evidence{07_pipeline_chunks}{Ejecución con chunks, mensaje STOP y suma correcta.}

Rank 0 crea un arreglo de $1\,048\,576$ enteros. El tamaño de chunk debe
dividir al total exactamente; por ello todos los bloques tienen la misma
cantidad de elementos. El productor conserva hasta dos envíos
\texttt{MPI\_Isend} pendientes. Antes de reutilizar una petición, espera
su finalización. Rank 1 detecta el siguiente mensaje, recibe el chunk y
procesa sus enteros de inmediato. Al terminar todos los envíos, rank 0
envía mediante \texttt{MPI\_Isend} un mensaje de control \texttt{STOP}.
El receptor lo reconoce por su etiqueta y sale del bucle.

Cada entero se compara con el patrón esperado y se acumula una suma.
Al final, el consumidor comunica al productor el número de chunks y su
suma. Rank 0 compara ambos contra los valores previstos. El tiempo total
se toma en rank 0 y termina después de la confirmación del consumidor y
una barrera final; incluye transferencia y procesamiento del lote completo.

La figura textual muestra varios chunks en el trayecto productor-consumidor
y un segundo mensaje de control STOP. La función \texttt{produce}
administra las dos peticiones pendientes; \texttt{consume} inspecciona
etiqueta y tamaño de cada mensaje, procesa el bloque y sale del bucle
únicamente tras recibir STOP. La figura se genera después de la medición.

El número de chunks es $C=M/B$, donde $M$ es la cantidad total de enteros
y $B$ el tamaño de chunk en enteros. Chunks pequeños implican más mensajes
y posiblemente mayor costo de coordinación; chunks grandes reducen el
número de mensajes, pero pueden limitar el solapamiento. La tendencia
concreta debe decidirse con los datos observados.

\clearpage
\subsubsection*{Inciso 3: tiempo total frente al tamaño de chunk}
\evidence{08_pipeline_tamanos}{Pipeline con otro tamaño de chunk: volumen fijo y distinto número de mensajes.}

\subsubsection*{Resultados del barrido}
El volumen total permanece en 1\,048\,576 enteros de cuatro bytes:
4\,194\,304 bytes, o 4 MiB. Los tamaños son divisores exactos de ese
total y la ventana conserva como máximo dos peticiones de envío.

\begin{center}\small
\begin{tabular}{@{}rrrrr@{}}
\toprule
\textbf{Enteros/chunk} & \textbf{Chunks} & \textbf{Mediana (ms)} & \textbf{Media $\pm$ DE (ms)} & \textbf{Mín.--máx. (ms)}\\
\midrule
256 & 4096 & 4.691 & $4.413\pm0.848$ & 3.462--5.088\\
1024 & 1024 & 3.670 & $3.667\pm0.097$ & 3.569--3.762\\
4096 & 256 & 2.844 & $2.940\pm0.333$ & 2.666--3.311\\
16384 & 64 & 2.769 & $3.589\pm1.824$ & 2.318--5.679\\
65536 & 16 & 2.279 & $2.271\pm0.118$ & 2.149--2.385\\
\bottomrule
\end{tabular}\end{center}

\clearpage
\chart{04_pipeline_chunks}{Tiempo total frente a bytes por chunk. Se mantiene un volumen de 4 MiB; las barras muestran el rango de las tres corridas.}

\subsubsection*{Tamaño de bloque, dispersión y corrección}
El menor tiempo mediano corresponde a chunks de 65\,536 enteros
(262\,144 bytes), con 2.279 ms. Respecto a los chunks de 256 enteros,
el tiempo mediano se reduce 51.42\,\% y la relación entre duraciones
es 2.059. Es una comparación entre tamaños dentro del pipeline,
no una aceleración frente a ejecución secuencial. Se identifica el
mejor tamaño ensayado, no un óptimo universal.

Los mensajes de datos disminuyen de 4096 a 16 sin alterar el volumen.
El comportamiento es compatible con una reducción de la sobrecarga
por mensaje; no se cronometraron por separado la transferencia, el
procesamiento y la planificación para atribuir causalidad a una sola etapa.

El caso de 16\,384 enteros presenta el mayor rango, 2.318--5.679 ms.
Su media de 3.589 ms supera a la mediana de 2.769 ms. Se conserva la
corrida de 5.679 ms sin diagnosticar una causa no medida. Las medianas
decrecen con el tamaño, mientras las medias no son monótonas; por eso
se muestran las observaciones individuales además del resumen.

Las 15 sumas son 523\,641\,600, exactamente la suma de $i\bmod1000$
para el arreglo completo. El consumidor compara cada elemento, suma
en orden y confirma la cantidad de bloques. La señal STOP se completa
antes de liberar buffers; las capturas muestran la terminación del
bucle y las verificaciones correctas.

\clearpage
\heading{Discusión metodológica}

Los experimentos sobre una sola máquina y WSL pueden reflejar costos de
memoria compartida, planificación del sistema operativo y configuración
de Open MPI, además de los costos del protocolo. El uso de
\texttt{-{}-oversubscribe} en Token Ring también puede introducir espera
por CPU si se exceden los recursos; la opción por sí sola no demuestra
esa condición. Los tiempos entre ejercicios no se comparan como si
midieran un mismo recurso. Cada gráfica responde a la pregunta concreta
del ejercicio y señala su escala y unidad.

Las salidas verificadas son una condición previa para interpretar los
tiempos. En el conjunto recibido no hay filas fallidas ni incompletas y
no se descartó ningún resultado. Tres repeticiones permiten describir
un mínimo, un máximo y una tendencia central, pero proporcionan poca
evidencia sobre estabilidad a largo plazo. Los tamaños se recorren en
orden fijo, no aleatorizado; cambios de carga o temperatura durante el
barrido podrían influir. No existe una traza de esas condiciones.

\subheading{Qué mide y qué no mide cada experimento}
Ping-Pong mide un ciclo completo y promedia 1000 rondas. El anillo mide
el protocolo de circulación hasta la barrera final. La recepción
anticipada mezcla espera, cálculo y consultas locales. El pipeline
incluye además procesamiento, comprobación de cada elemento, STOP y
confirmación. No deben interpretarse sus diferencias como una
clasificación global de primitivas MPI.

El contador de iteraciones no es una tasa de FLOP. La presencia de
\texttt{Isend} o \texttt{Irecv} no demuestra por sí misma transferencia
asíncrona autónoma ni una aceleración: describe la interfaz no bloqueante
y la posibilidad de continuar trabajo antes de completar la petición.
La mejora entre tamaños del pipeline no es \emph{speedup} frente a un
programa secuencial. Estas distinciones evitan conclusiones que los datos
no respaldan.

\subheading{Reproducibilidad y futuras extensiones}
Los cuatro CSV se versionan junto con código, capturas y scripts. La
reproducción del análisis no requiere nuevas ejecuciones MPI. Las
figuras se exportan en PDF vectorial y PNG; el informe incorpora la
versión vectorial. Los registros completos del barrido se conservan
localmente fuera de Git y se comprobaron contra los CSV.

Una ampliación del experimento podría aumentar repeticiones, aleatorizar
el orden, controlar afinidad y carga, y añadir una referencia bloqueante
con trabajo idéntico para estudiar el solapamiento. También convendría
ensayar más tamaños de chunk y máquinas conectadas por red. Estas
posibilidades no se presentan como pruebas ya realizadas.

\clearpage
\heading{Conclusiones}

El diseño cubre los cuatro ejercicios del enunciado con parámetros
controlables, comprobaciones de corrección y un protocolo explícito de
medición. La comunicación bloqueante exige coordinar el orden de envíos y
recepciones; \texttt{MPI\_Sendrecv} simplifica el desplazamiento circular.
Las operaciones no bloqueantes permiten iniciar comunicación y continuar
el cálculo o la producción de datos, siempre que se controle la vida de
los buffers y la finalización de las peticiones.

\begin{enumerate}
\item Ping-Pong conserva el último mensaje y exhibe crecimiento del RTT
mediano de 0.495 a 50.171 $\mu$s. El costo no es proporcional únicamente
al número de bytes y la mitad del RTT sigue siendo una estimación.
\item Ambas variantes del anillo devuelven token 5000 tras 1000 vueltas
con cinco ranks. La variante separada registra 2.242 ms frente a 3.290 ms
con Sendrecv, pero los 5000 y 25\,000 mensajes respectivos impiden
atribuir toda la diferencia al costo de la primitiva.
\item La recepción anticipada verifica todas las sumas y realiza trabajo
antes de confirmar la recepción. Sus medianas van de 6.701 a 121.223
$\mu$s; estos tiempos incluyen cálculo y sondeos y no cuantifican una
aceleración frente a una referencia bloqueante.
\item El pipeline procesa 4 MiB y conserva la suma 523\,641\,600 en
las 15 corridas. El mejor tamaño ensayado es 65\,536 enteros por chunk,
con 2.279 ms y una reducción mediana de 51.42\,\% frente al bloque mínimo.
\item Las 51 observaciones se mantienen. La mayor dispersión del pipeline
en chunks de 16\,384 enteros hace necesario mostrar rango, media y
mediana, en lugar de seleccionar solo la corrida más favorable.
\item La corrección del protocolo, la finalización de peticiones y la
vida de los buffers son condiciones previas para interpretar rendimiento.
Las conclusiones se limitan al equipo, configuraciones y muestra disponibles.
\end{enumerate}

\clearpage
\heading{Referencias}

\begin{enumerate}
\raggedright
\item Universidad del Valle de Guatemala. \emph{Laboratorio \#3},
Computación Paralela y Distribuida, semestre 2 de 2025. Enunciado PDF
proporcionado para esta realización.
\item Open MPI. \emph{MPI\_Sendrecv}. Manual oficial.
\url{https://docs.open-mpi.org/en/main/man-openmpi/man3/MPI_Sendrecv.3.html}.
\item Open MPI. \emph{MPI\_Test}. Manual oficial.
\url{https://docs.open-mpi.org/en/main/man-openmpi/man3/MPI_Test.3.html}.
\item Open MPI. \emph{MPI\_Wtime}. Manual oficial.
\url{https://docs.open-mpi.org/en/main/man-openmpi/man3/MPI_Wtime.3.html}.
\end{enumerate}

\clearpage
\heading{Apéndice: ejecución y archivos experimentales}

\begin{center}
\fcolorbox{azul}{azulclaro}{\parbox{0.90\linewidth}{%
\ttfamily\small
make\\
mpirun -np 2 ./build/ping\_pong 10 1\\
mpirun -np 2 ./build/ping\_pong 1000 1024 -{}-warmup\\
mpirun -{}-oversubscribe -np 5 ./build/token\_ring sendrecv 1\\
mpirun -np 2 ./build/recepcion\_anticipada 1 1000\\
mpirun -np 2 ./build/pipeline\_chunks 1048576 4096\\
bash scripts/mediciones.sh\\
python3 scripts/analizar\_resultados.py\\
python3 scripts/graficas\_resultados.py
}}
\end{center}

El barrido de mediciones genera cuatro CSV y conserva las salidas completas
de cada corrida. Se detiene si ya existen los CSV para evitar reemplazarlos;
un nuevo barrido requiere archivar previamente tanto tablas como registros.
Los dos scripts Python analizan los archivos existentes sin ejecutar MPI.
El de gráficas requiere Matplotlib; el analizador usa únicamente la
biblioteca estándar de Python.

\subheading{Archivos incorporados y trazabilidad}

\begin{center}
\begin{tabular}{@{}L{0.30\linewidth}L{0.58\linewidth}@{}}
\toprule
\textbf{Archivo} & \textbf{Uso en este informe} \\
\midrule
\texttt{ping\_pong.csv} & Gráfica de bytes frente a ida y vuelta media \\
\texttt{token\_ring.csv} & Comparación descriptiva de las dos variantes \\
\texttt{recepcion\_anticipada.csv} & Tiempo y trabajo efectuado según bytes \\
\texttt{pipeline\_chunks.csv} & Gráfica de bytes por chunk frente a tiempo total \\
Capturas PNG & Evidencia visible antes del análisis de cada apartado \\
\bottomrule
\end{tabular}
\end{center}

Cada CSV incluye tres repeticiones por configuración. Los registros
originales de terminal se conservan por separado, de modo que cualquier
valor de una figura pueda rastrearse hasta el comando y la ejecución que
lo produjo. Las ocho capturas recibidas se conservaron sin alterar su
contenido y se nombraron según las referencias de la fuente LaTeX.
Las cuatro gráficas proceden exclusivamente del barrido CSV, no de
transcribir las demostraciones de consola.

\subheading{Comprobación de los requisitos}
\begin{center}\small
\begin{tabular}{@{}L{0.31\linewidth}L{0.60\linewidth}@{}}
\toprule
\textbf{Requisito} & \textbf{Implementación y evidencia}\\
\midrule
Dos ranks, entero y N rondas & Ping-Pong parametrizado; capturas, tabla y gráfica de tamaños.\\
Anillo con más de cuatro ranks & Cinco ranks, vecinos modulares, dos variantes y verificación de visitas.\\
Recepción y aritmética con Test & Irecv publicado antes del cálculo; contador y estado visibles; gráfica de tamaños.\\
Chunks iguales y STOP con Isend & Divisibilidad validada, ventana de dos envíos, procesamiento inmediato y suma correcta.\\
Informe y código fuente & Fuente LaTeX, PDF, cuatro programas documentados, README, CSV y scripts; ejecutables excluidos.\\
\bottomrule
\end{tabular}\end{center}

\end{document}
